%%% FIELD proof-uniform-single-wave-risk
Fix a member of \(\mathfrak M\), a parameter \(\beta\in\mathcal K\), a balanced set \(T\) of size \(k\geq Q\), an action \(H\in\mathcal H_T\), and a realized bounded potential-outcome table.  Let \(U_T\) denote the uniform probability mass function on \(\Omega_T\), put \(w_q=2^{1-K_f}c_q\), and set
\[
L=\eta+\frac{1-\eta}{\alpha}.
\]
By \(\mathrm{def{:}exact\text{-}mixture}\), \(0\leq a_{T,H}(z)\leq1\), and hence
\[
\frac{p_{T,H}(z)}{U_T(z)}
=\eta+\frac{1-\eta}{\alpha}a_{T,H}(z)
\leq L.
\tag{1}
\]
The positive-support component of the same mixture and the balanced uniform marginal \(U_T(Z_i=q)=1/Q\) give, for every \(i\in T\) and \(q\in\mathcal Q\),
\[
\pi_{iq}(T,H)
=\sum_{z\in\Omega_T}p_{T,H}(z)\mathbf 1\{z_i=q\}
\geq \eta\sum_{z\in\Omega_T}U_T(z)\mathbf 1\{z_i=q\}
=\frac{\eta}{Q}>0.
\tag{2}
\]
Thus every division by an inclusion probability below is valid.

The exact-inclusion calculation underlying \(\mathrm{thm{:}exact\text{-}ht\text{-}accounting}\) also applies to this one wave.  Conditional on the realized table,
\[
\begin{aligned}
E_{p_{T,H}}\!\left\{\widehat{\boldsymbol\tau}^{\mathrm{HT}}_T(H)\mid\mathbf Y_n,\mathbf X_n\right\}
&=\frac1k\sum_{i\in T}\sum_{q\in\mathcal Q}
w_qY_{in}(q)\frac{P_{p_{T,H}}(Z_i=q)}{\pi_{iq}(T,H)}\\
&=\frac1k\sum_{i\in T}\sum_{q\in\mathcal Q}w_qY_{in}(q)
=\boldsymbol\tau_T.
\end{aligned}
\tag{3}
\]
Here the last equality is the definition of \(\boldsymbol\tau_T\).  Notice that (3) uses the actual inclusion probabilities, so neither unequal inclusions nor boundary randomization changes the centering.

Let \(e\) be a vector in an orthonormal basis of \(\mathbb R^F\), and define
\[
f_i(q)=\frac{e^{\mathsf T}W^{1/2}w_qY_{in}(q)}{\pi_{iq}(T,H)},
\qquad
X_e=\frac1k\sum_{i\in T}f_i(Z_i),
\qquad
c_e=E_{U_T}X_e.
\tag{4}
\]
Under balanced complete randomization, direct counting gives, for distinct \(i,j\in T\),
\[
P_{U_T}(Z_i=q,Z_j=r)
=\frac{(k/Q)(k/Q-\mathbf 1\{q=r\})}{k(k-1)}.
\tag{5}
\]
Substitution of (5) into the covariance definition yields
\[
\operatorname{Cov}_{U_T}\{f_i(Z_i),f_j(Z_j)\}
=-\frac1{k-1}\operatorname{Cov}_{q\sim\operatorname{Unif}(\mathcal Q)}
\{f_i(q),f_j(q)\}.
\tag{6}
\]
Consequently, expanding the variance of the sum and then collecting the covariances over the common uniform treatment label gives
\[
\begin{aligned}
\operatorname{Var}_{U_T}\!\left\{\sum_{i\in T}f_i(Z_i)\right\}
&=\frac{k}{k-1}\sum_{i\in T}\operatorname{Var}_{q}\{f_i(q)\}
-\frac1{k-1}\operatorname{Var}_{q}\!\left\{\sum_{i\in T}f_i(q)\right\}\\
&\leq \frac{k}{k-1}\sum_{i\in T}E_q\{f_i(q)^2\}.
\end{aligned}
\tag{7}
\]
The final term on the first line is nonpositive because it is minus a variance.

By (3), \(E_{p_{T,H}}X_e=e^{\mathsf T}W^{1/2}\boldsymbol\tau_T\).  Variance minimizes mean squared distance from a constant, and (1) compares the two assignment laws.  Therefore
\[
\begin{aligned}
E_{p_{T,H}}\!\left[
\left\{X_e-e^{\mathsf T}W^{1/2}\boldsymbol\tau_T\right\}^2
\mid\mathbf Y_n,\mathbf X_n\right]
&=\operatorname{Var}_{p_{T,H}}(X_e)\\
&\leq E_{p_{T,H}}\{(X_e-c_e)^2\}\\
&\leq L E_{U_T}\{(X_e-c_e)^2\}\\
&=L\operatorname{Var}_{U_T}(X_e).
\end{aligned}
\tag{8}
\]
The model-class condition \(\mathrm{ass{:}factor\text{-}count}\) gives \(Q=2^{K_f}\geq4\); because \(k\geq Q\), this implies \(k/(k-1)\leq4/3\).  Combining (2), (4), (7), and (8) thus gives
\[
\begin{aligned}
\operatorname{Var}_{p_{T,H}}(X_e)
&\leq \frac{L}{k^2}\frac{k}{k-1}
\sum_{i\in T}\frac1Q\sum_{q\in\mathcal Q}
\frac{\{e^{\mathsf T}W^{1/2}w_qY_{in}(q)\}^2}{\pi_{iq}(T,H)^2}\\
&\leq \frac{4LQ}{3\eta^2k^2}
\sum_{i\in T}\sum_{q\in\mathcal Q}
\{e^{\mathsf T}W^{1/2}w_qY_{in}(q)\}^2.
\end{aligned}
\tag{9}
\]
Summing (9) over an orthonormal basis and using
\[
\sum_e\{e^{\mathsf T}W^{1/2}w_qY_{in}(q)\}^2
=Y_{in}(q)^2w_q^{\mathsf T}Ww_q
\tag{10}
\]
proves the conditional bound
\[
\begin{aligned}
&E_{p_{T,H}}\!\left[
(\widehat{\boldsymbol\tau}^{\mathrm{HT}}_T(H)-\boldsymbol\tau_T)^{\mathsf T}
W(\widehat{\boldsymbol\tau}^{\mathrm{HT}}_T(H)-\boldsymbol\tau_T)
\mid\mathbf Y_n,\mathbf X_n\right]\\
&\qquad\leq
\frac{4LQ}{3\eta^2k^2}
\left\{\max_{q\in\mathcal Q}w_q^{\mathsf T}Ww_q\right\}
\sum_{i\in T}\sum_{q\in\mathcal Q}Y_{in}(q)^2.
\end{aligned}
\tag{11}
\]
The bounded-outcome condition ensures that the terms are integrable for each member of \(\mathfrak M\).  Averaging (11) under \(P_{\beta,n}(d\mathbf Y_n\mid\mathbf X_n)\), as prescribed by the definition of \(V_T\), proves the displayed inequality in the lemma.

For the consequence, suppose a subclass has a constant \(C<\infty\) such that
\[
\sup\frac1k\sum_{i\in T}\sum_{q\in\mathcal Q}
E_\beta\{Y_{in}(q)^2\mid x_{in}\}\leq C,
\tag{12}
\]
where the supremum is over its members, parameters, and balanced subsets.  Substitution of (12) into the proved bound yields \(V_T(H,\beta)\leq C'/k\), with \(C'\) independent of the member, \(\beta\), \(T\), and \(H\).  This is the asserted uniform \(O(k^{-1})\) conclusion.

Finally, the qualification concerning the whole class is substantive.  Given any admissible bounded regular member with nonzero design risk, transform every potential outcome by \(Y(q)\mapsto M Y(q)\) for a finite constant \(M>0\).  This preserves memberwise boundedness, smoothness, the score information under the one-to-one sample-coordinate transformation, and the normalized oracle action, while multiplying every quadratic design risk by \(M^2\).  Since \(M\) is not bounded across members by any primitive condition defining \(\mathfrak M\), those primitives cannot supply a single numerical risk constant uniform over all of \(\mathfrak M\).  Hence the lemma correctly asserts uniform first-order risk only on subclasses satisfying (12).
